\section{Introduction}
\label{sec:introduction}

Hall-algebra realizations of quantum cluster algebras are strongest when they
explain both the individual cluster variables and the multiplication of
arbitrary elements.  For a finite acyclic quiver, Ding--Xu--Zhang construct a
subquotient of the Hall algebra of morphisms between projectives which realizes
the principal-coefficient quantum cluster algebra
\cite{DingXuZhang}.  Fu--Peng--Zhang recover that construction from a bialgebra
and an integration map and identify the quantum Caldero--Chapoton formula
\cite{FuPengZhang}.  Those theorems use the hereditary setting in an essential
way.  The present project asks what survives when the algebra is a
nonhereditary cluster-tilted algebra.

Let $A$ be finite-dimensional and
\[
 \MM=\Mor(\proj A)
\]
the exact category of maps between finitely generated projective right
$A$-modules.  Bautista describes its indecomposable objects, and Li proves that
its maximal rigid objects are controlled by support $\tau$-tilting pairs
\cite{Bautista,Li}.  More precisely, the noncontractible indecomposable objects
have the form
\[
 C_M=(P_1^M\longrightarrow P_0^M),
 \qquad
 Z_P=(P\longrightarrow0),
\]
while $K_P=(P\xrightarrow{1}P)$ is contractible.  Li's correspondence sends a
support $\tau$-tilting pair $(M,P)$ to
\[
 C_M\oplus Z_P\oplus K_A
\]
and identifies the maximal-rigid exchange graph with the support
$\tau$-tilting graph.  Thus the category already contains the correct rigid
combinatorics beyond the hereditary case.

The preceding iterations of this work separated three layers.
\begin{enumerate}[label=(\roman*),leftmargin=2.3em]
\item The complete homology pair of a projective morphism reconstructs its
      reduced two-term object and its split index.
\item On reachable rigid objects, that information reconstructs the classical
      cluster character, the quantum seed and the quantum theta label.
\item The ordinary exact Hall integration exists for $\MM$, but its lattice
      exponent remembers only the two projective terms and forgets the
      differential.  A global bicharacter twist corrects the multiplication on
      every rigid compatibility cone.
\end{enumerate}
The third point leaves a sharp obstruction.  If $X$ and $Y$ are compatible,
then $\Ext^1(X,Y)=\Ext^1(Y,X)=0$ and their Hall product is a normalized direct
sum.  The corrected Hall product therefore agrees with the based product of
the corresponding quantum cluster monomials.  If $X$ and $Y$ are
incompatible, the Hall product contains non-split extensions, often with
non-rigid middle terms.  The union of the rigid chamber algebras is not closed,
and an object-level theta assignment is not yet an algebra homomorphism.

This paper resolves that last layer for cluster-tilted algebras of Dynkin type
$A$.  The restriction is mathematically substantial.  In type $A_n$, the
cluster category has the convex $(n+3)$-gon model: indecomposable objects are
diagonals, compatibility is non-crossing, and exchange triangles are the two
smoothings of a crossing.  The associated cluster-tilted algebras are gentle,
and extensions of string modules are described geometrically by smoothing
crossings \cite{CanakciSchroll}.  On the quantum side, the localized skein
algebra of an unpunctured marked surface is a quantum cluster algebra
\cite{Muller}; coefficient systems can be incorporated by walled skein
algebras \cite{IshibashiKanoYuasa}, and the recent stated-skein theorem proves
the equality with the frozen-localized quantum cluster and upper cluster
algebras for polygons, together with the theta-basis identification
\cite{CaoHuangWang}.

\subsection*{The Hall--skein ideal}

Fix a triangulation $T$ of the polygon $P=P_{n+3}$ and the corresponding
cluster-tilted algebra $A$.  Work over the coefficient field
$\mathbb Q(\omega)$, with $q$ specialized to the prescribed power of
$\omega$.  Starting with the exact Hall algebra of $\MM$, we perform the same
three operations already isolated in the earlier paper:
\begin{enumerate}[label=(\alph*),leftmargin=2.3em]
\item localize at the contractible classes $[K_P]$;
\item apply the Hall-compatible term-lattice bicharacter twist;
\item centrally specialize contractible directions to the frozen boundary
      monomials of the chosen coefficient system.
\end{enumerate}
The resulting algebra is denoted $\Hred$.  In type $A$, every indecomposable
reduced object is rigid and corresponds to one polygon diagonal.  Since the
category is representation-finite, the usual triangular Hall argument shows
that these indecomposable classes generate all of $\Hred$.

For every ordered crossing pair $\gamma,\delta$, the quantum skein algebra has
a local relation
\begin{equation}
 s_\gamma s_\delta
 =\omega^{a(\gamma,\delta)}s_{\gamma\#_0\delta}
  +\omega^{b(\gamma,\delta)}s_{\gamma\#_\infty\delta}.
 \label{eq:intro-skein}
\end{equation}
Here each smoothing is a two-component multicurve, and a boundary component is
an invertible frozen variable.  We impose the same relation on the
corresponding Hall generators and let $\Jsk$ be the generated two-sided ideal.
Our main theorem states
\begin{equation}
 \boxed{
  \Hred/\Jsk
  \;\xrightarrow{\sim}\;
  \Skein_\omega^{\rm fr}(P)
  \;\xrightarrow{\sim}\;
  \Aq.
 }
 \label{eq:intro-main}
\end{equation}
This is a genuine subquotient realization of the whole quantum cluster
algebra.  The source is obtained from the entire exact Hall algebra, not from
a vector space spanned only by rigid objects.  In particular, arbitrary Hall
products are defined before quotienting, and the quotient assigns a quantum
cluster element to every non-rigid middle term.

\subsection*{Why the quotient is not tautological}

One could always define an ideal as the kernel of a desired map.  That is not
what is done here.  The ideal $\Jsk$ is given explicitly by the local crossing
relations \eqref{eq:intro-skein}.  The proof that no further relations are
needed uses the crossing filtration.

For a reduced object $M$, let $\crs(M)$ be the number of crossings in the
multidiagonal formed by its Krull--Schmidt summands.  The split middle term in
an ordered Hall product is the unique term of maximal crossing number.  By the
extension--smoothing theorem, every non-split middle term has strictly smaller
crossing number.  Hence a Hall--skein defect has an invertible leading
coefficient times the crossing direct sum.  It may therefore be oriented as a
reduction of one crossing basis vector to lower-crossing basis vectors.

The remaining issue is confluence.  We prove a filtered Diamond lemma.  A
critical ambiguity is a multidiagonal with two possible crossings to resolve.
The difference of the two reductions is a lower-crossing element of the same
Hall--skein ideal.  Induction reduces it to zero; geometrically this is the
usual local confluence of polygon skein smoothing, and algebraically it uses
the associativity of the Hall product.  Thus the irreducible classes are
precisely the non-crossing multidiagonals.  Muller's skein basis theorem shows
that the same set is a basis of the localized polygon skein algebra, so the
surjection in \eqref{eq:intro-main} is injective.

\subsection*{Non-rigid middle terms}

The quotient provides an explicit answer to the question which motivated the
present iteration.  Let $M$ be any Hall basis object.  If $M$ is non-rigid,
choose a crossing pair among its indecomposable summands.  Expand the actual
Hall product of that pair, solve for the split direct-sum term, and replace the
product by the two skein smoothings.  All remaining Hall middle terms have
smaller crossing number.  Repeating the procedure gives a finite expression
\begin{equation}
 \St_{\mathrm H}(M)
 =\sum_{R\text{ non-crossing}}
    c_{M,R}(\omega)\,\eps_R.
 \label{eq:intro-straightening}
\end{equation}
The Diamond lemma makes the answer independent of all choices.  Under
\eqref{eq:intro-main}, $\eps_R$ is the normalized quantum cluster monomial and,
equivalently, the theta function $\Theta_{g(R)}$.  Therefore every non-rigid
middle term obtains the canonical theta expansion
\[
 \Phi_{\rm sk}(\eps_M)
 =\sum_R c_{M,R}(\omega)\Theta_{g(R)}.
\]
No positivity claim is made for the Hall-corrected coefficients
$c_{M,R}(\omega)$; the theorem is an algebraic realization and a normal-form
statement.

For arbitrary $x,y\in\Hred$, including incompatible rigid generators, the
quotient homomorphism satisfies
\[
 \Phi_{\rm sk}(x\star y)
 =\Phi_{\rm sk}(x)\Phi_{\rm sk}(y).
\]
Thus the Hall structure constants of every incompatible product are converted
into quantum/theta structure constants after the explicit skein quotient.
This is the type-$A$ analogue of the global subquotient phenomenon of
Ding--Xu--Zhang and Fu--Peng--Zhang.

\subsection*{Scope and relation to existing multiplication formulas}

Caldero--Keller realize a finite-type cluster algebra as an exceptional Hall
algebra of its cluster category \cite{CalderoKeller}; Palu proves cluster
multiplication formulas for arbitrary objects in suitable $2$-Calabi--Yau
categories \cite{PaluMultiplication}.  Weighted quantum and motivic refinements
for arbitrary pairs are developed by Chen--Xiao--Xu and
Xiao--Xu--Yang \cite{ChenXiaoXu,XiaoXuYang}.  These results supply the
$2$-Calabi--Yau multiplication background.  Our construction has a different
source and conclusion: it starts from the exact Hall algebra of projective
morphisms over a nonhereditary cluster-tilted algebra and gives an explicit
local ideal whose quotient is the quantum cluster algebra.

The theorem is not claimed for every Jacobi-finite cluster-tilted algebra.
Outside the polygon and related surface models, an explicit skein presentation
for all non-rigid objects is unavailable, and the literature explicitly notes
that a general quantum Laurent expansion for arbitrary $2$-Calabi--Yau
objects is not known \cite{XiaoXuYang}.  The final section isolates the exact
input needed for a future generalization: a terminating and confluent
extension-straightening system whose standard classes are the theta basis.

\subsection*{Organization}

\Cref{sec:prior} recalls the morphism-category and Hall framework.
\Cref{sec:skein} fixes the polygon skein algebra and coefficients.
\Cref{sec:crossing} proves the crossing-filtration and Diamond-lemma results.
\Cref{sec:subquotient} proves the global subquotient theorem.
\Cref{sec:nonrigid} gives the non-rigid theta straightening algorithm.
\Cref{sec:A4} records the complete type-$A_4$ certificate.
\Cref{sec:audit} separates direct proofs, imported theorems and remaining
questions, and \cref{sec:reproduction} describes the code and reproducibility
package.
